\section{Conclusion}
\label{sec:conclusion}

We present the first direct measurement of how steering vector influence decays in transformer hidden states after the vector is removed. Using teacher-forced generation to control for token identity, we show that the steering signal drops to less than 5\% of its peak within a single token, following an exponential decay with time constants $\timeconstant \approx 1.2$--$5.8$ tokens. The KV cache is the primary vehicle for any residual steering signal: preserving it retains a small positive trace, while resetting it causes a statistically significant overcorrection (Cohen's $d$ up to 3.41). These findings confirm that continuous steering is necessary for sustained behavioral control via activation addition, and provide mechanistic support for KV-cache-based methods as a more persistent alternative.

Several directions for future work emerge. First, repeating this analysis across model scales (7B, 70B+) would test whether decay rate is scale-dependent. Second, comparing decay across different steered behaviors (\eg refusal, language, reasoning style) would reveal whether some behaviors produce more persistent perturbations. Third, the KV cache overcorrection phenomenon warrants further investigation: does it have downstream behavioral consequences, and can it be exploited for steering in the opposite direction? Finally, measuring decay under multi-layer steering would test whether distributing the perturbation across layers increases persistence.
